% Created 2026-09-05 Sat 13:25
% Intended LaTeX compiler: xelatex
\documentclass[12pt]{article}
\usepackage[utf8]{inputenc}
\usepackage[T1]{fontenc}
\usepackage{graphicx}
\usepackage{longtable}
\usepackage{wrapfig}
\usepackage[margin=1in]{geometry}
\usepackage{rotating}
\usepackage{fontspec}
\setmainfont{Libertinus Serif}
\usepackage{setspace}
\onehalfspacing
\usepackage[normalem]{ulem}
\usepackage{amsmath}
\usepackage{parskip}
\usepackage{amssymb}
\usepackage{capt-of}
\setcounter{secnumdepth}{4} % Tells LaTeX to number down to paragraphs
\setcounter{tocdepth}{4}
\usepackage{hyperref}
\usepackage[style=oscola, backend=biber]{biblatex}
\addbibresource{/Users/krishnajani/Documents/Libib.bib}
\usepackage{csquotes}
\usepackage[british]{babel}
\setlength{\parskip}{0.5\baselineskip}
\author{Krishna Jani}
\date{\today}
\title{National Exploration Policy - Hydrocarbons Exploration Policy}

\begin{document}

\maketitle

\begin{enumerate}
\item Before 1991-92, petroleum exploration and production was carried out only by the public sector oil companies such as ONGC and Oil India Limited (OIL). A New Exploration and Licensing Policy (NELP) was formulated in 1997 for exploration and production of crude oil and natural gas (excluding CBM - Coal Based Methane).
\item The main objective was to attract significant risk capital from Indian and foreign companies, state-of-the-art technologies, geological concepts, and best international practices for exploration and production of oil and gas resources in the country
\item A Director General of Hydrocarbons established in 1993, was appointed as the nodal agency for the new NELP policy
\item Another significant boost to this sector was when exploration and production was opened up to private and foreign investment with allowance of 100\% FDI under the rubric of the NELP scheme
\item The following are the key features of NELP
    \begin{enumerate}
        \item FDI of upto 100\% is allowed under NELP
        \item Blocks are to be awarded through an open international competitive bidding mechanism.

        NELP operated on periodic discrete bidding rounds (NELP-1 through NELP-19) and companies had to wait for the Ministry to officially release a set of predefined blocks during a specific window.
        \item ONGC and OIL are at the same level as private companies and have to compete for obtaining the petroleum exploration licenses on a competitive basis rather than obtaining them on nomination basis as it existed prior. They also get the same fiscal and contractual terms as private companies
        \item The contractors are allowed to freely market their crude oil and gas in the domestic market, although export of oil and gas is not permitted
        \item Two-tiered royalty system where 12.5\% for onland areas and 10\% for offshore areas to incentivize the exploration of hydrocarbons offshore. This extends to twelve nautical miles from the coastline of India.
        \item There is no mandatory state participation through ONGC or OIL, neither is there any carried interest of the state.

        Normally under a carried interest setup the government or the state-owned company gets a free percentage share in an oil block. The private company spends 100\% of the capital investment required for drilling, survey and exploration and if no oil is found the private company loses its money while the government is owed nothing. Conversely if oil is found the private company covers the upfront cost and gradually recovers the expense out of the future oil sales once it reaches a profit making stage it is split between the state and the private entity.

        The objective here is to allow the state to take a share in the successful project without taking up high upfront financial risk. Before the NELP ONGC and OIL were granted oil fields on a nomination basis and offers held mandatory stakes or carried rights in the blocks allowing them a split of the profit without the excessive upfront financial overhead. In order to create a level playing field the automatic free equity stakes granted to ONGC and OIL were no longer valid
        \item Royalty was charged at half the prevailing rate for companies that engaged in deep water mining beyond 400 meters for the first seven years after commencement. This was because mining at a depth of 400 meters and above is harder and more expensive than drilling near the shore. Therefore, to enable companies to recover their huge setup costs and overheads, an exception was made.
        \item Cess was exempted for production from blocks offered under NELP.
        \item Companies were exempted from payments of import duty in respect of goods imported for petroleum operations
        \item No signature discovery or production bonuses. Meaning that companies need not pay extra upfront or performance based cash gifts to the government such as a payment upon signing a contract or upon discovery or bonus paid once the oil and gas production reaches a specific volume target.

        Focus is again to remove the upfront financial burden on private and foreign companies, making it attractive to explore the hydrocarbon sector in India
    \end{enumerate}
\end{enumerate}
\section*{Issues with NELP}
There are separate policy regimes for conventional oil and gas, coal bed methane, shale oil and gas hydrates. Different fiscal terms are also in force for allocation of acreages and different hydrocarbons. However, in practice there is an overlapping of resources between different contracts. Unconventional hydrocarbons such as shale gas and shale oil, which refer to natural gas and liquid crude oil trapped under organic formations called ``shale'', were unknown when NELP was framed.

The fragmented policy framework led to inefficiencies where upon exploring one type of hydrocarbon if another is found it will need separate licensing which adds to the overhead cost

NELP was also based on PSC that is production sharing contracts which are based on the principle of profit sharing when a contractor discovers oil and gas, he is expected to share with the government the profit from his venture until a profit is made no share is given to the government. Now since the contract requires profit to be measured, it also becomes necessary for the costs to be checked and verified by the government. For this the government had to create a system of approval at various stages to prevent the contractor from exaggerating the costs and activities could not be commenced until this approval was given. This gave the government a lot of discretion and became a major source of delay and dispute since there was sharp disagreement between the government and the contractor leading to regarding the necessity or lack of necessity of particular costs or the correctness of correctness of the cost. There was also revenue leakage in certain cases with the certain corporations showing inordinate losses on their production and exploration process while in fact this these were mere devices to evade profit sharing

Exploration was confined to blocks that were put to tender by the government. However, private entities might also have certain blocks that they would like to pursue for exploration based on their own information or interest, but this was an untapped opportunity till the time the government brought that those blocks for bidding

The pricing of gas, especially the producer price of gas was fixed administratively by the government. In most cases the government seeks the price at lower than the actual market value, which led to lower corporate revenue and therefore lower taxes and royalties. Although it is stated that there was freedom to price oil and gas in the domestic market, this freedom was subject to a price ceiling imposed by the government. In cases where the global market was soaring, the price ceiling ensured that the production companies could not take advantage of the soaring prices while in cases where the global prices were experiencing a crash the applicable price was below the ceiling.

\end{document}
